\section{Utilidades}
\subsection{Subset sum optimization.}
Complejidad: $O(N\sqrt{N})$. Si queremos obtener calcular el conjunto de sumas posibles en subconjuntos de un arreglo $A$, en el cual se cumple que $\sum a_i = N$ y todo $a_i\geq0$, podemos obtener exactamente el mismo conjunto con un arreglo comprimido de tamaño $\sqrt{N}$. El $3k-trick$ explota que las sumas generadas por el arreglo $[a,a,a]$ son exactamente las mismas que en $[2a,a]$. Podemos juntar elementos y repetirlo hasta que a solo haya dos copias de cada elemento distinto en el nuevo arreglo, acotando el numero de elementos a $O(\sqrt{N})$. El proceso de compresión toma $O(N\log_2(N))$.
\begin{minted}{cpp}
vi compress(vector<ll>& a){
    sort(a.rbegin(), a.rend());
    int n = sz(a);
    a.pb(0);
    vi weights;
    int pi = 0;
    for(int i = 1; i <= n; i++){
      if(a[i] != a[i - 1]){
        ll cnt = i - pi;
        ll x = a[i - 1];
        ll j = 1;
        while(j < cnt){
          weights.pb(x * j);
          cnt -= j;
          j *= 2;
        }
        weights.pb(x * cnt);
        pi = i;
      }
    } return weights;
}
\end{minted}



\subsection{Bitsets de tamaño (casi) dinámico}
El código ejecutado dentro de func(sz) usará un bitset con tamaño de la potencia de 2 más cercano a sz.
\begin{minted}{cpp}
template<int len = 1> void func(int sz){
    if(sz > len){
      func<std::min(len * 2, MAXLEN)>(sz);
      return;
    }
    bitset<len> bs;
    //usa aquí tu bitset dinámico
}
\end{minted}



\subsection{Bitwise ops}
\begin{minted}{cpp}
#define is_on(S, j) (S & (1ll << (j)))
#define set_bit(S, j) (S |= (1ll << (j)))
#define clear_bit(S, j) (S &= ˜(1ll << (j)))
#define toggle_bit(S, j) (S ˆ= (1ll << (j)))
#define lsb(S) ((S) & -(S))
#define clear_lsb(S) (S &= (S - 1))
#define set_all(S, n) (S = (1ll << (n)) - 1ll)
#define clear_trailing_ones(S) (S &= (S + 1))
#define set_last_bit_off(S) (S |= (S + 1))
#define is_power_of_two(S) (!((S) & ((S) - 1)))
#define nearest_power_of_two(S) ((int)pow(2, (int)((log((double)(S)) / log(2)) + 0.5)) )
#define is_divisible_by_power_of_two(n, k) !((n) & ((1ll << (k)) - 1))
#define modulo(S, N) ((S) & ((N) - 1)) // S%N, N=2^k
\end{minted}

\vspace{-0.7\baselineskip}

\subsection{Bitwise - builtin}
\begin{minted}{cpp}
// one plus the index of the least significant 1-bit of x, or if x is zero, returns zero.
int __builtin_ffs(int x)
// number of leading 0-bits in x, starting at the most significant bit position. If x is 0 is undefined
int __builtin_clz(unsigned int x)
// number of trailing 0-bits in x, starting at the least significant bit position. If x is 0  undefined.
int __builtin_ctz(unsigned int x)
// he parity of x, i.e. the number of 1-bits in x modulo 2.
int __builtin_parity(unsigned int x)
\end{minted}



\subsection{Iterar}
Dada una máscara m, iterar sobre todos sus subconjuntos desde el más grande hasta el más pequeño (ojo el código no itera sobre $x=m$). La complejidad de iterar sobre todas las submáscaras de todos los números de 1 a $2^n $es $O(3^n)$.
\begin{minted}{cpp}
void submasks(int mask){
    for( int x = mask; x;){
        --x &= mask;
    }
}
\end{minted}



\subsection{Gospers’ Hack}
Sirve para generar todos las máscaras de $n$ bits, que tengan exactamente $k$ bits a $1$ (y que sean menores o iguales que $2^n$). \\
Complejidad $O\pa{\binom{n}{k}}$?
\begin{minted}{cpp}
void GospersHack(int n, int k) {
    int mask = (1 << k) - 1;
    int limit = (1 << n);
    while(mask < limit){
        int c = mask & - mask;
        int r = mask + c;
        mask = (((r ^ mask) >> 2) / c) | r;
    }
}
\end{minted}



\subsection{Subset Sum con bitset}
Dado una una arreglo de tamaño \(n\) ver si es psoible encontrar es la suma objetivo  \(S\). Nota: \(w\) es el tamaño del bloque de bits procesados en paralelo (normalmente 32 o 64). Complejidad:  Temporal \(O\left(n \times \frac{S}{w}\right)\) -  Memoria \(O(S)\)

\begin{minted}{cpp}
bool subsetSumBitset( vi& arr, int S) {
    bitset<10001> dp; // tam >= S+1
    dp[0] = 1;
    for(int x : arr) dp |= (dp << x);
    return dp[S];
}
\end{minted}



\subsection{SOS DP}

Dado un conjunto de $n$ elementos y una función $A(S)$ definida para cada subconjunto $S$, queremos calcular
\[
F(S) = \sum_{T \subseteq S} A(S), \quad \forall S \subseteq \{0,1,\dots,n-1\}.
\]

Definimos 
$
S(x, i) = \{ y \subseteq x \mid y \oplus x < 2^i \},
$ 
notar que contiene todos los subconjuntos de $x$ cuyos bits a la izquierda del bit $i$ coinciden con los de $x$.

\textbf{Ejemplo:}  
$$
S(11010110_2, 3) = \{ 
  11010000, 
  11010010, 
  11010100, 
  11010110 \}
$

Podemos descomponer $S(x, i)$ recursivamente:

\[
S(x, i) =
\begin{cases}
S(x, i-1) & \text{si el bit $i$ de $x$ es 0} \\
S(x, i-1) \cup S(x \oplus 2^i, i-1) & \text{si el bit $i$ de $x$ es 1}
\end{cases}
\]
Con esta partición, definimos la tabla de DP:
\[
dp[x][i] = \sum_{y \in S(x, i)} f[y]
\]
Y asi tenemos
\[
dp[x][i] =
\begin{cases}
dp[x][i-1] & \text{si el bit $i$ de $x$ es 0} \\
dp[x][i-1] + dp[x \oplus 2^i][ i-1] & \text{si el bit $i$ de $x$ es 1}
\end{cases}
\]


\textbf{T:} $O(n \cdot 2^n)$ | \textbf{M:} $O(2^n)$.
\begin{minted}{cpp}
int n;
vector<int> f(1<<n); // inicializar f[S]
vector<int> g = f;

for(int i = 0; i < n; i++) {
    for(int mask = 0; mask < (1<<n); mask++) {
        if(mask & (1<<i)) {
            g[mask] += g[mask^(1<<i)];
        }
    }
}
\end{minted}
